\documentclass[aspectratio=169]{beamer}

%-----------------------------------
% Tema e pacotes
\usetheme{Madrid}
\usecolortheme{default}

\usepackage[utf8]{inputenc}
\usepackage{graphicx}
\usepackage{booktabs}
\usepackage{amsmath}

%-----------------------------------
% Dados do título
\title[Tema 7]{Tema 7 – Procura Sistemática de Literatura}
\author{Departamento de Física \\ Universidade Eduardo Mondlane}
\date{\today}

%-----------------------------------
\begin{document}

%-----------------------------------
\begin{frame}
\titlepage
\end{frame}

%-----------------------------------
\begin{frame}{Conteúdo}
\tableofcontents
\end{frame}

%-----------------------------------
\section{Introdução}

\begin{frame}{Introdução}
\begin{itemize}
    \item Revisões sistemáticas consolidam evidências científicas.
    \item Avaliam estudos de forma objetiva, estruturada e transparente.
    \item A pesquisa sistemática garante rigor metodológico.
\end{itemize}
\end{frame}

%-----------------------------------
\begin{frame}{Objetivos da Sessão}
\begin{itemize}
    \item Compreender a importância da pesquisa sistemática.
    \item Desenvolver estratégias de pesquisa robustas.
    \item Aplicar frameworks como PICO/ST.
    \item Explorar técnicas avançadas de pesquisa.
\end{itemize}
\end{frame}

%-----------------------------------
\section{Revisão Sistemática}

\begin{frame}{Principais Características}
\begin{itemize}
    \item Estruturada, objetiva e reproduzível.
    \item Critérios de elegibilidade bem definidos.
    \item Estratégias de pesquisa abrangentes.
    \item Síntese rigorosa dos resultados.
\end{itemize}
\end{frame}

%-----------------------------------
\begin{frame}{Características da Revisão Sistemática}
\centering
\includegraphics[width=0.75\textwidth]{caracteristicas_rsm.png}
\end{frame}

%-----------------------------------
\section{Pesquisa Sistemática de Literatura}

\begin{frame}{Importância da Pesquisa Sistemática}
\begin{itemize}
    \item Identificação abrangente de estudos relevantes.
    \item Redução de viés de seleção.
    \item Base para sínteses qualitativas e quantitativas.
\end{itemize}
\end{frame}

%-----------------------------------
\begin{frame}{Por que usar pesquisa sistemática?}
\centering
\includegraphics[width=0.75\textwidth]{importancia_pesquisa.png}
\end{frame}

%-----------------------------------
\section{Estratégia de Pesquisa}

\begin{frame}{Etapas do Processo}
\begin{enumerate}
    \item Definição da questão de pesquisa (PICO/ST)
    \item Identificação de palavras-chave
    \item Uso de operadores booleanos
    \item Pesquisa em bases de dados
    \item Aplicação de critérios de elegibilidade
    \item Organização dos resultados
\end{enumerate}
\end{frame}

%-----------------------------------
\begin{frame}{Fluxo da Pesquisa Sistemática}
\centering
\includegraphics[width=0.75\textwidth]{etapas_pesquisa.png}
\end{frame}

%-----------------------------------
\section{Framework PICO/ST}

\begin{frame}{Framework PICO/ST}
\begin{itemize}
    \item \textbf{P} – População
    \item \textbf{I} – Intervenção / Exposição
    \item \textbf{C} – Comparador
    \item \textbf{O} – Desfecho
    \item \textbf{S} – Tipo de estudo
    \item \textbf{T} – Tempo
\end{itemize}
\end{frame}

%-----------------------------------
\begin{frame}{Exemplo de Aplicação PICO/ST}
\begin{table}
\centering
\begin{tabular}{p{2cm} p{7cm}}
\toprule
Elemento & Exemplo \\
\midrule
P & Pacientes com diabetes tipo 2 \\
I & Dieta de baixo índice glicêmico \\
C & Dieta padrão \\
O & HbA1c \\
S & Ensaios clínicos randomizados \\
T & 6 meses \\
\bottomrule
\end{tabular}
\end{table}
\end{frame}

%-----------------------------------
\section{Técnicas de Pesquisa}

\begin{frame}{Palavras-chave e Sinônimos}
\begin{itemize}
    \item Identificação de termos principais.
    \item Uso de sinônimos e variações.
    \item Tesauros das bases de dados.
\end{itemize}
\end{frame}

%-----------------------------------
\begin{frame}{Operadores Booleanos}
\begin{itemize}
    \item \textbf{OR}: sinônimos
    \item \textbf{AND}: conceitos diferentes
    \item \textbf{NOT}: exclusão de termos
\end{itemize}

\centering
\includegraphics[width=0.6\textwidth]{operadores_booleanos.png}
\end{frame}

%-----------------------------------
\begin{frame}{Truncamento e Proximidade}
\begin{itemize}
    \item Truncamento: uso do * para variações.
    \item Proximidade: termos próximos no texto.
\end{itemize}

\centering
\includegraphics[width=0.6\textwidth]{truncamento_proximidade.png}
\end{frame}

%-----------------------------------
\section{Bases de Dados}

\begin{frame}{Execução da Pesquisa}
\begin{itemize}
    \item Scopus
    \item Web of Science
    \item PubMed
    \item Google Scholar
    \item Uso de filtros (ano, idioma, tipo de estudo)
\end{itemize}
\end{frame}

%-----------------------------------
\begin{frame}{Bases de Dados}
\centering
\includegraphics[width=0.75\textwidth]{bases_dados.png}
\end{frame}

%-----------------------------------
\section{Exemplo Prático}

\begin{frame}{Exemplo de Consultas}
\begin{table}
\centering
\begin{tabular}{p{3cm} p{7cm}}
\toprule
Base & Consulta \\
\midrule
Scopus & TITLE-ABS-KEY(diabetes AND "low glycemic index") \\
WoS & TS=(diabetes AND "low glycemic index") \\
PubMed & ("diabetes"[Mesh] AND "low glycemic index") \\
\bottomrule
\end{tabular}
\end{table}
\end{frame}

%-----------------------------------
\section{Exercício}

\begin{frame}{Exercício Prático}
\begin{itemize}
    \item Formular questão PICO/ST.
    \item Criar estratégia de pesquisa.
    \item Executar em uma base de dados.
    \item Organizar os resultados.
\end{itemize}
\end{frame}

%-----------------------------------
\begin{frame}{Atividade}
\centering
\includegraphics[width=0.75\textwidth]{exercicio_pesquisa.png}
\end{frame}

%-----------------------------------
\section*{Referências}

\begin{frame}{Referências}
\small
\begin{itemize}
    \item Higgins et al. \textit{Cochrane Handbook}, 2019.
    \item Borenstein et al. \textit{Introduction to Meta-Analysis}, 2009.
    \item Moher et al. PRISMA Statement, 2009.
\end{itemize}
\end{frame}

%-----------------------------------
\begin{frame}
\centering
\Huge Obrigado!
\end{frame}

\end{document}
